\section{Related Work}
\label{sec:related}

Our research is situated at the intersection of cooperative perception, multi-agent motion forecasting, and the system-level design of latency-aware edge platforms. Prior work has advanced these areas individually but has not addressed how delay and reordering can make the edge world model internally inconsistent, or how those inconsistencies affect closed-loop safety.

\subsection{Cooperative Perception Architectures}
The goal of cooperative perception is to enable vehicles to share sensor data to overcome the line-of-sight limitations of individual agents. These systems vary in their architecture and the level of data they share.

\textbf{Infrastructure-less V2V systems}, such as AutoCast~\cite{Autocast}, exchange sensor data to improve local perception but require high-density deployments and reliable, low-latency links. They also remain limited by line-of-sight: if no vehicle in the ad-hoc network can see a hidden hazard, that information cannot be shared. An edge-based system with a roadside sensor on a high vantage point can overcome this limitation. AutoCast scales dense settings by scheduling which vehicles share at each step, so it does not directly test the high-participation failure mode we study.

\textbf{Edge-assisted V2I systems} address the shortcomings of infra\-structure-less V2V systems by leveraging roadside infrastructure for more robust data fusion and intelligence generation. Early work like EMP~\cite{EMP} demonstrated the feasibility of fusing raw LiDAR point clouds at an edge server. More recent systems have focused on optimizing different parts of the pipeline. For example, VIPS~\cite{vips2022} compensates for communication latency between vehicles and infrastructure by temporally aligning stale detections using velocity-based extrapolation at the edge. We implement a faithful version of the VIPS temporal alignment mechanism as one of our three evaluation baselines to specifically isolate and study its limits under dynamic maneuvers. EdgeCooper~\cite{edgecooper2023} focuses on optimizing the V2X communication layer itself by proposing a scheduling algorithm to intelligently select which data to share to maximize perception utility under bandwidth constraints. Harbor~\cite{harbor} further explored edge-assisted object fusion.
%\ad{Cite this more recent work? Tyler: This is covered in the paragraph below (V2X-VITl, OpenCOOD, etc}

\noindent\textbf{Identity anchoring vs. feature extrapolation.}
Learned fusion and staleness-mitigation (e.g., extrapolating features or tracks under latency) improve accuracy but remain probabilistic. They cannot prevent duplicate ego births when messages arrive out of order. In contrast, our edge-enforced identity anchoring provides a deterministic invariant (one ego identity $\rightarrow$ one track) and eliminates duplicates by construction, independent of motion model or feature drift.


\begin{figure}[t]
    \centering
    % Placeholder for the high-level system overview diagram.
    \includegraphics[width=\columnwidth]{figures/Vehicle Architecture - Page 1-5.pdf}
    \caption{High-level overview of our edge-assisted perception system. Vehicles and RSUs perform local perception and transmit lightweight 3D object lists and self-beacons to the edge. The edge node runs our intelligence pipeline to create and disseminate a unified, predictive world model.}
    \label{fig:system_overview}
\end{figure}

%\ad{you should make figures denser (less whitespace). this can fit in a single column if packed properly. Packed tighter, moved to single column}

\noindent\textbf{Benchmarks and evaluation scale.}
Most cooperative-perception benchmarks evaluate a small number of collaborating agents per scene and report perception metrics (e.g., AP) rather than closed-loop safety.
For example, OPV2V~\cite{xu2022opencood} constrains each frame to 2--7 collaborating vehicles, DAIR-V2X~\cite{DAIR-V2X} focuses on V2I perception with explicit temporal asynchrony handling, and V2X-Sim~\cite{V2xsim} selects 2--5 vehicles per scene as collaboration agents. These benchmarks are valuable for model comparison, but they do not chart closed-loop safety envelopes as the number of concurrently participating agents and tail staleness grow.


\subsection{Motion Forecasting and Planning}

Effective edge intelligence requires perceiving the present and predicting the future.  Prior research in motion forecasting has produced powerful models for this purpose.

\textbf{Multi-modal Trajectory Prediction} is critical for handling the inherent uncertainty of traffic scenarios. Influential works like Trajectron++~\cite{trajectron2020} and MultiPath++~\cite{multipath2021} use deep learning to generate multiple plausible future trajectories for road agents. Other research has focused on specific challenges, such as producing uncertainty-aware predictions~\cite{djuric2020uncertainty} or handling complex interactions in specific domains like roundabouts~\cite{li20203}. These models, however, are typically evaluated on offline datasets and assume perfect, zero-latency perception data. Our work measures what happens when predictors consume the noisy, stale data produced by a real-world edge pipeline.

\textbf{Cooperative Motion Prediction and Planning} takes motion prediction a step further by considering explicit communication between agents. For example, CMP~\cite{cmp_ieee_2023} explores how multi-agent communication can improve motion predictions. These approaches assume reliable, low-latency communication channels. Our work analyzes the failure modes that arise when those assumptions are violated.

%\subsection{Learned vs. Model-Based Motion Prediction}
%While our work uses a  simple Kalman filter, we acknowledge powerful learned predictors (e.g., LSTMs, SSMs) such as Mamba~\cite{huang2025mambamotstatespacemodelmotion, mambatrack} can model complex, non-linear motion. However, these models introduce a trade-off between predictive accuracy and verifiability, acting as `black boxes' that are difficult to certify for safety-critical systems like ours (ISO 26262)~\cite{Serna2020}. Furthermore, their higher computational cost would consume more of the tight edge compute budget, potentially exacerbating the very latency-induced failures we aim to solve~\cite{fusionlatencywacv2025}. We therefore argue that for this problem, our deterministic, protocol-level solution is more robust, verifiable, and efficient than relying on a more complex prediction model.

\subsection{Positioning Our Contribution}
Prior cooperative perception work evaluates perception quality using proxy metrics such as Average Precision or processing latency. These metrics measure how well the edge sees the world. They do not measure whether the planner can safely act on what the edge publishes. We close this gap in two ways.

First, we evaluate safety outcomes directly: collision rates, false-braking events, and operational success under delay variability. Measuring closed-loop outcomes reveals failure modes such as self-ghosting that are invisible in proxy-metric evaluations.

Second, we design self-beacon anchoring, an edge-native protocol that enforces ego-uniqueness at the publish boundary. Prior systems such as EMP~\cite{EMP} and Harbor~\cite{harbor} focus on the accuracy of fused data. Anchoring addresses a different problem, preventing the edge from publishing the ego's own stale track as an obstacle.

\noindent\textbf{QoS vs. invariants.}
Retry, priority, and admission heuristics reduce average staleness but do not offer correctness guarantees under reordering. Without an identity invariant, proximity gates or hysteresis can still birth a duplicate ego track at moderate latency. We therefore treat identity as an application-level invariant rather than a QoS parameter.


%\KR{@Tyler: kudos! the related work section is comprehensive and well-written.}

